\DocumentMetadata{
  pdfversion=2.0,
  pdfstandard=ua-2,
  lang=en-US,
  tagging=on,
  tagging-setup={math/setup=mathml-SE,tabsorder=structure}
}
\documentclass[11pt,letterpaper]{article}
\usepackage[margin=0.68in,includefoot]{geometry}
\usepackage{newcomputermodern}
\usepackage{amsmath,amscd,mathtools}
\usepackage{tikz,adjustbox}
\providecommand{\square}{\mdlgwhtsquare}
\usepackage[hidelinks]{hyperref}
\hypersetup{
  pdftitle={Analysis First-Year Exam, May 2020, Part 2},
  pdfauthor={University of Florida Department of Mathematics},
  pdfsubject={Graduate mathematics examination},
  pdfdisplaydoctitle=true,
  bookmarksopen=true
}
\setcounter{secnumdepth}{0}
\setlength{\parindent}{0pt}
\setlength{\parskip}{0.28em}
\setlength{\emergencystretch}{3em}
\setlength{\leftmargini}{2.15em}
\setlength{\leftmarginii}{2.65em}
\setlength{\leftmarginiii}{3em}
\pagestyle{empty}
\raggedbottom
\allowdisplaybreaks
\NewTaggingSocketPlug{block/list/label}{examproblem}{%
  \tagstructbegin{tag=Lbl}\tagstructend%
  \tagstructbegin{tag=\LBody}%
  \tagstructbegin{tag=\ProblemHeadingTag,title-o={\ProblemHeadingText}}%
  \tagmcbegin{tag=\ProblemHeadingTag,actualtext={\ProblemHeadingText}}%
  #2%
  \tagmcend%
  \tagstructend%
}
\newenvironment{examproblems}[2]{%
  \def\ProblemHeadingTag{#2}%
  \AssignTaggingSocketPlug{block/list/label}{examproblem}%
  \begin{enumerate}%
  \setcounter{enumi}{#1}%
  \renewcommand{\labelenumi}{\textbf{\theenumi.}}%
  \setlength{\topsep}{0.35em}%
  \setlength{\partopsep}{0pt}%
  \setlength{\itemsep}{0.48em}%
  \setlength{\parsep}{0.18em}%
}{\end{enumerate}}
\newenvironment{examsubparts}{%
  \AssignTaggingSocketPlug{block/list/label}{default}%
  \begin{enumerate}%
  \setlength{\topsep}{0.18em}%
  \setlength{\partopsep}{0pt}%
  \setlength{\itemsep}{0.16em}%
  \setlength{\parsep}{0.1em}%
}{\end{enumerate}}
\newenvironment{examinstructions}{%
  \par\begingroup\itshape
}{%
  \par\endgroup\vspace{0.18em}
}
\makeatletter
\renewcommand\subsection{\@startsection{subsection}{2}{0pt}%
  {-1.25ex plus -.25ex minus -.1ex}%
  {0.55ex plus .1ex}%
  {\normalfont\large\bfseries}}
\renewcommand{\@maketitle}{%
  \newpage\null\vskip -1em%
  {\footnotesize\bfseries
    \href{https://www.ufl.edu/}{University of Florida}, \href{https://math.ufl.edu/}{Department of Mathematics}\par}%
  \vskip -0.5em%
  \tagstructbegin{tag=H1,title={Analysis First-Year Exam, May 2020, Part 2}}%
  \tagpdfparaOff%
  \tagmcbegin{tag=H1}%
    {\LARGE\bfseries\href{https://gma.math.ufl.edu/exams/analysis-first-year/}{Analysis First-Year Exam}, May 2020, Part 2\par}%
  \tagmcend%
  \tagpdfparaOn%
  \tagstructend%
  \vskip 0.25em%
  \tagmcbegin{artifact}%
    \hrule height 1pt%
  \tagmcend%
  \vskip 1em%
}
\makeatother
\title{Analysis First-Year Exam, May 2020, Part 2}
\author{}
\date{}
\begin{document}
\maketitle
\thispagestyle{empty}
\begin{examinstructions}
Answer FOUR questions in detail. State carefully any results used without proof.
\end{examinstructions}
\begin{examproblems}{0}{H2}
\def\ProblemHeadingText{Problem 1.}
\item
Let \(a<b\) and let \(f:[a,b]\to\mathbb{R}\) be Riemann integrable. Prove that if \(\int_a^b f=0\), then for each positive integer~\(n\) there exist \(c<d\) in \([a,b]\) such that \(\sup\{f(t):c\leq t\leq d\}<1/n\). Hence, or otherwise, prove that if \(f>0\) throughout \([a,b]\), then \(\int_a^b f>0\).
\def\ProblemHeadingText{Problem 2.}
\item
Let \((x_n:n\geq 0)\) converge in the metric space~\(X\) and let \((f_n:n\geq 0)\) be a sequence of continuous real-valued functions on~\(X\).
\begin{examsubparts}
\item[\textnormal{(a)}]
Prove that if the sequence of functions is uniformly convergent, then \((f_n(x_n):n\geq 0)\) is convergent.
\item[\textnormal{(b)}]
Show by example that if uniformly is replaced by pointwise, then \((f_n(x_n):n\geq 0)\) may not converge.
\end{examsubparts}
\def\ProblemHeadingText{Problem 3.}
\item
Let \(f:[0,1]\to\mathbb{R}\) be continuous and assume that
\[
\int_0^1 f(t)t^n\,dt=1/(n+1)
\]
 for each integer \(n>1\). Deduce as much as is possible about~\(f\).
\def\ProblemHeadingText{Problem 4.}
\item
Let \((f_n:n\geq 0)\) be a sequence of measurable real-valued functions on some measurable space. Prove that each of the following sets is measurable:
\[
A=\{\omega:f_n(\omega)\in[0,1]\text{ for finitely many }n\}
\]

\[
B=\{\omega:f_n(\omega)\in[0,1]\text{ for infinitely many }n\}
\]

\[
C=\left\{\omega:\text{the series }\sum_{n=0}^{\infty}f_n(\omega)\text{ is absolutely convergent}\right\}.
\]
\def\ProblemHeadingText{Problem 5.}
\item
For each positive integer~\(n\), let \(f_n:[0,1]\to[0,1]\) be continuous; assume that \(f_n\to 0\) pointwise as \(n\to\infty\). Does it follow that
\[
\int_0^1 f_n(t)\,dt\to 0\quad\text{as }n\to\infty?
\]
 Does your answer change if continuous is replaced by Riemann integrable?
\end{examproblems}
\end{document}
